\documentclass[a4paper, 12pt]{article}
\usepackage[margin=2cm]{geometry}
\usepackage{amssymb, amsmath, graphicx, tabularx, booktabs, subfiles, amsthm,enumitem,xcolor,appendix}
 

\newcommand{\ans}[1]{\textcolor{red}{#1}}

\newcommand{\boxx}[1]{%
    \begin{tcolorbox}[colback=gray!10, colframe=gray!50, title=Where it is getting wrong?]
        #1
    \end{tcolorbox}%
}


\title{}
\author{}

\begin{document}
\maketitle
%------------Body-Starts--------------------
% Text book document
\section{Book Text}
One author says: “Don’t serve more than two foods rich in sugar or starch at the same meal. When you serve bread and potatoes, your starchlicense has run out. A meal that includes peas, bread, potatoes, sugar, cake and after dinner mints should also include a Vitamin B Complex capsule, some bicarbonate of soda (other than that used on the vegetables), and the address of the nearest specialist in arthritis and other degenerative diseases.” For more than seventy years it has been the rule in Hygienic circles to take but one starch at a meal and to consume no sweet foods with the starch meal. Sugars, syrups, honeys, cakes, pies, mints, etc., have been tabu with starches. We do not say to those who come to us for advice: If you eat these with your starches, take a dose of baking soda with them. We tell them to avoid the fermentation that is almost inevitable. In Hygienic circles it is considered the height of folly to take a poison and then take an antidote with it. We think it best not to take the poison. Sugar with starch means fermentation. It means a sour stomach. It means discomfort. Those who are addicted to the honey-eating practice and who are laboring under the popular fallacy that honey is a “natural sweet” and may be eaten indiscriminately, should know that this rule not to take sweets with starches applies to honey as well. Honey or syrup, it makes no difference which, with your cereals, honey or sugar to sweeten your cakes – these combinations spell fermentation. White sugar, brown sugar, “raw” sugar, imitation brown sugar (that is, white sugar that has been colored), black strap molasses, or other syrup, with starches means fermentation. Soda will neutralize the resulting acids, it will not stop the fermentation. For more than sixty years it has been the practice in Hygienic circles to take a large raw vegetable salad (leaving out tomatoes or other acid foods) with the starch meal.

%My crisp understanding in about 10 lines
\subsection*{Understanding}
\begin{enumerate}
    \item \textbf{What is the author trying to say?} The author is trying to say that eating sugar and starch together is not good for health. It causes fermentation and sour stomach.
    \item \textbf{What is the author's opinion?} The author's opinion is that eating sugar and starch together is not good for health.
    \item \textbf{What is the author's evidence?} The author's evidence is that it has been the rule in Hygienic circles to take but one starch at a meal and to consume no sweet foods with the starch meal. Sugars, syrups, honeys, cakes, pies, mints, etc., have been tabu with starches.
    \item \textbf{What is the author's conclusion?} The author's conclusion is that eating sugar and starch together is not good for health.
\end{enumerate}
%Key Message in 3 lines
\subsection*{Key Message}
Eating sugar and starch together is not good for health. It causes fermentation and sour stomach. It has been the rule in Hygienic circles to take but one starch at a meal and to consume no sweet foods with the starch meal. Sugars, syrups, honeys, cakes, pies, mints, etc., have been tabu with starches.
%My crisp understanding in about 1 lines










%------------Body-Ends--------------------
%\bibliography{ref.bib}
%\bibliographystyle{IEEEtran}
\end{document}
